\documentclass[11pt]{article}
\usepackage[a4paper, top=18mm, bottom=18mm, left=20mm, right=20mm]{geometry}
\usepackage[T2A]{fontenc}
\usepackage[utf8]{inputenc}
\usepackage[ukrainian]{babel}
\usepackage{graphicx}
\usepackage[export]{adjustbox}
\usepackage{amsmath}
\usepackage{amssymb}
\graphicspath{{/app/Quiz/Scripts/2023/ProgAll/15BinTree/}{/app/Quiz/Scripts/2021/Mod2021_3/BinTree/}{/app/Quiz/Scripts/2020/Mod2020_4/BinTree/}}
\setlength{\parindent}{0pt}
\setlength{\parskip}{6pt}
\pagestyle{empty}
\begin{document}
%begin
{\large\textbf{Контрольна робота}}

\textbf{Варіант \No\,9}\hfill \textit{ПІБ:}\,\underline{\hspace{60mm}}
\vspace{4mm}

%beginitem
1. Що буде виведено за виконання фрагменту коду?

%code
print('R', end=' ')
i=6
while i<12:
    i+=1
print(i)
%code

\vspace{5mm}
%enditem

%beginitem
2. Що буде виведено за виконання фрагменту коду?

%code
i=6
while i<12:
    i+=1
print(i)
%code

\vspace{5mm}
%enditem

%beginitem
3. %goodpath:Scripts/2018/Mod2018_1/03-good.txt
Відмітити все, що є коректним оператором С++:
( 1 ) bool f(int a);

( 2 ) int f(int a,b);

( 3 ) cout<<3.4;

( 4 ) int n; bool n;
\vspace{5mm}
%enditem

%beginitem
4. Що буде виведено за виконання фрагменту коду?
code
_a = 0

class A:
    _a = 1
    
    def __init__(self):
        self._a = 2
        _a = 3
        A._a = 4

obj = A()
try:
    print(_a, end=' ')
    print(obj._a, end=' ')
    print(A._a, end=' ')
    print(obj._a, end=' ')
    print(obj._A__a, end=' ')
except:
    print('ERR')
code
\vspace{5mm}
%enditem

%beginitem
5. %goodpath:Scripts/2019/Mod2019_1/Lex/03-good.txt
Відмітити все, що є коректним оператором С++:
( 1 ) 3.14%2

( 2 ) if a<b: a=0 else: b=0

( 3 ) i,j=j,j

( 4 ) x+3=5
\vspace{5mm}
%enditem

%beginitem
6. Що буде виведено за виконання фрагменту коду?

%code
print('R', end=' ')
g=[10,11,12,13,14]
del g[6:]
print(g)
%code

\vspace{5mm}
%enditem

%beginitem
7. %goodpath:Scripts/2019/Mod2019_1/Lex/01-good.txt
Відмітити все, що є однією правильною лексемою:
( 1 ) False

( 2 ) '\'

( 3 ) py

( 4 ) 8,3
\vspace{5mm}
%enditem

%beginitem
8. Що буде виведено за виконання фрагменту коду?

%code
print('R', end=' ')
i=12
while i>6:
    i+=-3
print(i)
%code

\vspace{5mm}
%enditem

%beginitem
9. Що буде виведено за виконання фрагменту коду?
%code
if (15 < 15)
    cout << "x";
else
    cout << "c";
%code

\vspace{5mm}
%enditem

%beginitem
10. Що буде виведено за виконання фрагменту коду?

%code
def h(a,b):
    c=58
    if a<2:
        return 6
    if a>5:
         c=2
    else: 
        return 7
    return c

print(h(-2,-2))
%code
\vspace{5mm}
%enditem

%end
\newpage
\end{document}

